% id: 203-20
% book: 베이직쎈 대수 · 12 수학적 귀납법
% section: 유형 126 수학적 귀납법; 등식의 증명
% figure: none
% answer: 
% answer_source: 
% variant_level: 0 (원본 전사)
% 이 파일은 build.mjs 가 items.json 에서 만든다. 고칠 때는 items.json 을 고친다.
\smash{\raisebox{1.5mm}{\dmkichul{베이직쎈 대수 203쪽 20번}}}
다음은 모든 자연수 $n$에 대하여\par\smallskip\noindent\hspace*{1.5em}$1\times 2+2\times 3+3\times 4+\cdots+n(n+1)$\\ \hspace*{1.5em}$=\dfrac{n(n+1)(n+2)}{3}$\par\smallskip\noindent 가 성립함을 수학적 귀납법으로 증명한 것이다.\exprbox{\begin{minipage}{0.96\linewidth}\raggedright\lineskip=4pt {\small\bfseries 증명}\par\smallskip (i) $n=1$일 때,\\ \hspace*{3em}(좌변)$=1\times 2=2$, (우변)$=\dfrac{1\times 2\times 3}{3}=2$\\ \hspace*{1.5em}따라서 주어진 등식이 성립한다.\\ (ii) $n=k$일 때 주어진 등식이 성립한다고 가정하면\\ \hspace*{3em}$1\times 2+2\times 3+3\times 4+\cdots+k(k+1)$\\ \hspace*{1.5em}$=\dfrac{k(k+1)(k+2)}{3}$\\ \hspace*{1.5em}위의 식의 양변에 $\fbox{㈎}$를 더하면\\ \hspace*{3em}$1\times 2+2\times 3+3\times 4+\cdots+k(k+1)+\fbox{㈎}$\\ \hspace*{1.5em}$=\dfrac{k(k+1)(k+2)}{3}+\fbox{㈎}$\\ \hspace*{1.5em}$=\dfrac{\fbox{㈏}}{3}$\\ \hspace*{1.5em}따라서 $n=k+1$일 때도 주어진 등식이 성립한다.\\ (i), (ii)에서 모든 자연수 $n$에 대하여 주어진 등식이 성립한다.\end{minipage}}위의 과정에서 ㈎, ㈏에 알맞은 것을 각각 $f(k)$, $g(k)$라 할 때, $f(3)+g(2)$의 값은?
\dmchoicesauto{}{$80$}{$100$}{$120$}{$140$}{$160$}
